\section{Số liệu đánh giá}
\subsection{Các chỉ số đánh giá (Evaluation Metrics)}

Phần này trình bày các chỉ số đánh giá được sử dụng trong toàn bộ hệ thống, bao phủ từ mức tín hiệu (speech processing) đến mức suy luận ngữ nghĩa (RAG). Các chỉ số được thiết kế nhằm phản ánh đầy đủ các tiêu chí: độ chính xác, khả năng bảo toàn thông tin, và hiệu quả truy xuất.

\subsubsection{Diarization Metrics}

\textbf{Vai trò:}

Các chỉ số diarization được sử dụng để đánh giá chất lượng của mô-đun định danh người nói trong pipeline chuyển đổi âm thanh thành văn bản. Đây là bước nền tảng ảnh hưởng trực tiếp đến toàn bộ các thành phần phía sau như speech recognition, summarization, retrieval và question answering. Nếu speaker bị gán sai hoặc nội dung hội thoại bị mất trong giai đoạn diarization, các lỗi này sẽ tiếp tục lan truyền sang các tầng xử lý downstream và làm giảm chất lượng của toàn bộ hệ thống.

\textbf{Diarization Error Rate (DER):}

DER là chỉ số tiêu chuẩn được sử dụng để đo lường mức độ sai lệch giữa kết quả diarization và ground truth theo trục thời gian của hội thoại.

\[
DER = \frac{Miss + FalseAlarm + Confusion}{Total\ Speech\ Time}
\]

Trong đó:
\begin{itemize}
    \item \textbf{Miss (Missed Speech):} phần speech tồn tại trong audio nhưng không được hệ thống phát hiện.
    
    \item \textbf{False Alarm:} phần không phải speech nhưng bị nhận diện nhầm là speech.
    
    \item \textbf{Confusion:} phần speech được phát hiện nhưng bị gán sai speaker.
\end{itemize}

Trong phạm vi nghiên cứu này, quá trình phân tích tập trung chủ yếu vào hai thành phần Miss và Confusion. Miss làm mất thông tin đầu vào cho ASR và các bước xử lý phía sau, trong khi Confusion ảnh hưởng trực tiếp đến tính đúng đắn của speaker metadata. Điều này đặc biệt quan trọng đối với các hệ thống retrieval theo người tham gia hoặc các tác vụ tổng hợp thông tin theo speaker.

\textbf{Speaker Matching (Label Alignment):}

Do các speaker labels trong hệ thống diarization chỉ mang tính ký hiệu và không có ý nghĩa cố định, cần thực hiện bước ánh xạ nhãn giữa kết quả dự đoán và ground truth trước khi tính DER.

Bài toán này được mô hình hóa như một bài toán gán tối ưu và được giải bằng thuật toán Hungarian dựa trên confusion matrix ở mức frame. Quá trình alignment giúp đảm bảo việc tính toán DER phản ánh đúng chất lượng speaker assignment thay vì bị ảnh hưởng bởi sự khác biệt trong cách đánh số nhãn speaker.

\textbf{Phân tích lỗi:}

Ngoài các chỉ số định lượng, hệ thống còn thực hiện phân tích các dạng lỗi phổ biến trong speaker diarization nhằm hỗ trợ quá trình chẩn đoán và cải thiện pipeline.

\begin{itemize}
    \item \textbf{Label Oscillation:} hiện tượng speaker labels thay đổi liên tục trong cùng một đoạn phát biểu.
    
    \item \textbf{Over-merging:} nhiều speaker khác nhau bị gộp thành cùng một nhãn.
    
    \item \textbf{Over-splitting:} một speaker bị tách thành nhiều nhãn khác nhau.
\end{itemize}

Các phân tích này được sử dụng như công cụ hỗ trợ đánh giá hành vi của hệ thống và xác định nguyên nhân lỗi trong quá trình diarization, thay vì đóng vai trò như các chỉ số định lượng độc lập.

\textbf{Speaker-Aware Coverage (SAC):}

Speaker-Aware Coverage (SAC) là chỉ số được sử dụng để đánh giá đồng thời khả năng bảo toàn nội dung hội thoại và độ chính xác của speaker attribution trong transcript đầu ra. Khác với DER chủ yếu đánh giá lỗi diarization theo góc nhìn timeline và speaker segmentation, SAC tập trung vào chất lượng của transcript cuối cùng ở mức nội dung và ngữ nghĩa.

SAC được xây dựng từ hai thành phần chính gồm Coverage và Speaker Correctness.

\begin{itemize}
    \item \textbf{Coverage:} tỷ lệ các từ hoặc cụm từ mang ý nghĩa trong Ground Truth vẫn xuất hiện trong transcript đầu ra sau quá trình xử lý.
    
    \item \textbf{Speaker Correctness:} tỷ lệ nội dung được gán đúng speaker tương ứng.
\end{itemize}

Chỉ số SAC chỉ đạt giá trị cao khi hệ thống vừa giữ lại được nội dung hội thoại, vừa duy trì đúng mối liên hệ giữa nội dung và speaker metadata. Nếu transcript vẫn giữ được phần lớn nội dung nhưng speaker labels bị gán sai, hoặc ngược lại, SAC sẽ giảm đáng kể.

Điểm khác biệt quan trọng giữa DER và SAC nằm ở mục tiêu đánh giá. DER phản ánh lỗi diarization theo trục thời gian của audio, trong khi SAC phản ánh mức độ mà transcript cuối cùng vẫn bảo toàn đúng nội dung gắn với đúng speaker. Vì vậy, SAC phản ánh trực tiếp hơn tác động của lỗi diarization đối với các tác vụ downstream như retrieval theo participant, meeting summarization hoặc phân tích đóng góp của từng người tham gia trong cuộc họp.

\subsubsection{Summarization Metrics}

\textbf{Vai trò:}

Các chỉ số summarization được sử dụng để đánh giá khả năng của hệ thống trong việc rút gọn dữ liệu hội thoại dài thành biểu diễn cô đọng hơn nhưng vẫn giữ được các thông tin quan trọng phục vụ cho retrieval và indexing. Trong bối cảnh hệ thống tóm tắt đa tầng, việc đánh giá không chỉ tập trung vào độ ngắn gọn của bản tóm tắt mà còn cần xem xét khả năng bảo toàn nội dung và mức độ bám sát transcript gốc.

\textbf{Information Coverage Score (ICS):}

Information Coverage Score (ICS) được sử dụng để đo lường mức độ bảo toàn thông tin của hệ thống tóm tắt. Chỉ số này đánh giá xem các ý chính, quyết định, yêu cầu hoặc thông tin quan trọng trong ground truth có còn xuất hiện trong bản tóm tắt sinh ra hay không.

Quá trình đánh giá được thực hiện bằng cách phân tách nội dung thành các information units nhỏ, sau đó thực hiện semantic matching giữa ground truth và bản tóm tắt sinh ra thông qua embedding similarity. Hệ thống đồng thời kết hợp lexical overlap nhằm tăng độ ổn định của quá trình đánh giá coverage.

ICS phản ánh khả năng duy trì thông tin ở nhiều tầng summarization khác nhau như participant-level, topic-level và meeting-level. Giá trị ICS càng cao cho thấy bản tóm tắt càng giữ lại được nhiều nội dung quan trọng của cuộc họp.

\textbf{Compression Ratio (CR):}

Compression Ratio (CR) được sử dụng để đánh giá khả năng nén thông tin của hệ thống tóm tắt.

\[
CR = \frac{Summary\ Length}{Original\ Transcript\ Length}
\]

Chỉ số này phản ánh tỷ lệ giữa độ dài của bản tóm tắt đầu ra và transcript gốc. Giá trị CR thấp cho thấy hệ thống có khả năng rút gọn đáng kể kích thước dữ liệu cần lưu trữ và xử lý, đồng thời vẫn duy trì được lượng thông tin cần thiết phục vụ retrieval.

Trong bối cảnh hệ thống quản lý nhiều cuộc họp, khả năng giảm kích thước dữ liệu nhưng vẫn giữ được coverage là yếu tố quan trọng giúp giảm chi phí indexing và tăng tốc quá trình retrieval.
\textbf{Hallucination Rate (HR):}

Hallucination Rate (HR) được sử dụng để đánh giá mức độ xuất hiện của các thông tin không tồn tại trong transcript gốc.

Trong quá trình đánh giá, mỗi phát biểu trong bản tóm tắt sẽ được đối chiếu với transcript thông qua semantic similarity và lexical overlap. Nếu một phát biểu không đạt đủ mức tương đồng với nội dung gốc thì sẽ được xem là hallucination.

HR thấp cho thấy nội dung tóm tắt vẫn bám sát transcript ban đầu và hạn chế việc sinh ra các thông tin sai lệch hoặc không có thật. Điều này đặc biệt quan trọng khi bản tóm tắt được sử dụng như lớp trung gian phục vụ retrieval và indexing trong các hệ thống quản lý cuộc họp đa tài liệu.

\subsubsection{Các chỉ số đánh giá hệ thống RAG}

Để đánh giá toàn diện kiến trúc RAG, hệ thống sử dụng bộ chỉ số đo lường phân tách theo hai giai đoạn cốt lõi: Giai đoạn Truy xuất (Retrieval) và Giai đoạn Tạo sinh (Generation). Việc phân tách này giúp cô lập lỗi và xác định chính xác nguyên nhân làm giảm hiệu năng của hệ thống nằm ở khâu tìm kiếm ngữ cảnh hay khâu suy luận của Mô hình Ngôn ngữ Lớn (LLM).

\textbf{Đánh giá Mô-đun Truy xuất (Retrieval Metrics)}
Các chỉ số này tập trung đo lường độ chính xác của cơ sở dữ liệu vector và thuật toán tìm kiếm:

\begin{itemize}[leftmargin=*, label={}]
    \item \textbf{Hit Rate@K (Tỷ lệ trúng đích trong Top K):}
    Ý nghĩa: Đo lường tỷ lệ phần trăm các truy vấn mà hệ thống truy xuất thành công đoạn hội thoại chứa đáp án đúng nằm trong top K kết quả trả về. \\
    Mục đích: Kiểm chứng hiệu quả của chiến lược phân mảnh văn bản (\textit{Parent Document Retrieval}). So với phương pháp chia \textit{chunk} nguyên câu truyền thống, việc băm nhỏ câu thoại thành các đoạn 40 từ với độ chồng lấp 10 từ được kỳ vọng sẽ làm tăng đáng kể chỉ số này nhờ giảm thiểu tỷ lệ bỏ sót thông tin (High Recall).

    \item \textbf{Mean Reciprocal Rank - MRR (Điểm xếp hạng nghịch đảo trung bình):}
    \item  \[
    MRR = \frac{1}{Q} \sum_{i=1}^{Q} \frac{1}{rank_i}
    \]
    Ý nghĩa: Đánh giá chất lượng thứ hạng của đoạn văn bản chứa đáp án đúng (Ví dụ: đoạn văn bản đúng nằm ở Top 1 được 1 điểm, Top 2 được 0.5 điểm). \\
    Mục đích: Đảm bảo bộ tái xếp hạng (\textit{Cross-Encoder Re-ranker}) hoạt động hiệu quả, đẩy các ngữ cảnh quan trọng nhất lên đầu danh sách nhằm giảm tải dữ liệu "rác" (nhiễu) đưa vào LLM.

    \item \textbf{Context Relevance (Độ liên quan của Ngữ cảnh):}
    Ý nghĩa: Đo lường tỷ lệ giữa các thông tin thực sự hữu ích so với tổng lượng thông tin có trong cửa sổ ngữ cảnh (\textit{Context Window}). \\
    Mục đích: Khẳng định việc áp dụng kỹ thuật cửa sổ trượt (\textit{Sliding Window}) và mở rộng ngữ cảnh linh hoạt cung cấp vừa đủ lượng thông tin cần thiết cho LLM suy luận mà không làm loãng chủ đề truy vấn.
\end{itemize}

\textbf{Đánh giá Mô-đun Tạo sinh (Generation Metrics)}
Do đặc thù của dữ liệu âm thanh mang nhiều sai số (ASR Noise), việc đánh giá chất lượng câu trả lời cuối cùng do LLM sinh ra là bắt buộc nhằm kiểm chứng khả năng tự sửa lỗi và chống ảo giác:

\begin{itemize}[leftmargin=*, label={}]
    \item \textbf{Faithfulness (Tính xác thực / Độ bám sát ngữ cảnh):}
    Ý nghĩa: Đo lường xem câu trả lời của LLM có hoàn toàn dựa trên ngữ cảnh đã được truy xuất hay không, hay mô hình tự bịa đặt thêm thông tin. \\
    Mục đích: Đảm bảo tính trung thực của hệ thống. Một câu trả lời có thể đúng với kiến thức thực tế bên ngoài nhưng nếu không xuất phát từ biên bản cuộc họp thì vẫn bị đánh giá điểm Faithfulness thấp (vi phạm nguyên tắc \textit{grounded}).

    \item \textbf{Answer Relevance (Độ liên quan của câu trả lời):}
    Ý nghĩa: Đánh giá mức độ trực tiếp và trọn vẹn của câu trả lời so với câu hỏi gốc ban đầu. \\
    Mục đích: Đo lường khả năng lập luận đa bước (\textit{multi-step reasoning}) của hệ thống. Đối với các câu hỏi phức tạp, hệ thống phải tổng hợp đủ thông tin để đưa ra câu trả lời toàn diện thay vì trả lời vòng vo hoặc thiếu ý.
\end{itemize}

\textbf{Đánh giá Hiệu năng Siêu dữ liệu (Metadata Impact)}
Bên cạnh các chỉ số tiêu chuẩn, đồ án bổ sung thêm phương pháp đánh giá thực nghiệm (A/B Testing) nhằm đo lường vai trò của siêu dữ liệu:

\begin{itemize}[leftmargin=*, label={}]
    \item \textbf{Metadata Filtering Efficiency (Hiệu quả lọc siêu dữ liệu):}
    Ý nghĩa: So sánh chéo các chỉ số Hit Rate@K và MRR giữa hai cấu hình hệ thống: truy xuất vector thuần túy so với truy xuất vector kết hợp lọc siêu dữ liệu (định danh người nói, mốc thời gian). \\
    Mục đích: Chứng minh việc tổ chức siêu dữ liệu dạng \textit{payload} giúp loại bỏ triệt để các kết quả trùng lặp từ vựng nhưng khác nguồn phát biểu, đặc biệt quan trọng trong dữ liệu hội thoại đa diễn giả.
\end{itemize}